\documentclass[12pt]{article}
\usepackage{amsmath, amssymb}
\usepackage[margin=1in]{geometry}
\setlength{\parskip}{6pt}
\title{QUBO Formulation: Binary Support Vector Machine Candidate Selection}
\date{}
\begin{document}
\maketitle

\subsection*{Binary Support Vector Machine Candidate Selection}

\paragraph{Problem Statement.}
Given a set of $N$ candidate models, each associated with a precomputed regularized hinge-loss score $s_j$ for $j \in \{0, \dots, N-1\}$, the task is to select exactly one candidate such that the total score is minimized. The total score is defined as the sum of the classification loss and the regularization penalty, which are precomputed constants for each candidate.

\paragraph{Decision Variables.}
For each candidate $j \in \{0, \dots, N-1\}$, introduce a binary decision variable $x_j \in \{0, 1\}$. The variable $x_j$ acts as a selection indicator, where $x_j = 1$ denotes that candidate $j$ is chosen, and $x_j = 0$ otherwise. All variables are mutually exclusive selection indicators.

\paragraph{QUBO Derivation.}
\textbf{Step 1 --- Objective Function.} The original optimization problem is to minimize the weighted sum of candidate scores:
\begin{align}
  \min \quad \sum_{j=0}^{N-1} s_j x_j.
\end{align}

\textbf{Step 2 --- Constraint.} The requirement to select exactly one candidate is expressed as:
\begin{align}
  \sum_{j=0}^{N-1} x_j = 1.
\end{align}

\textbf{Step 3 --- Penalty Term Expansion.} To incorporate the constraint into the objective, we construct a quadratic penalty term with weight $\lambda > 0$:
\begin{align}
  P(\mathbf{x}) = \lambda \left( \sum_{j=0}^{N-1} x_j - 1 \right)^2.
\end{align}
Expanding the square yields:
\begin{align}
  P(\mathbf{x}) = \lambda \left( \sum_{j=0}^{N-1} x_j^2 - 2\sum_{j=0}^{N-1} x_j + 1 + \sum_{j \neq k} x_j x_k \right).
\end{align}
Applying the idempotent property $x_j^2 = x_j$ for binary variables simplifies the diagonal terms:
\begin{align}
  P(\mathbf{x}) = \lambda \left( \sum_{j=0}^{N-1} x_j - 2\sum_{j=0}^{N-1} x_j + 1 + \sum_{j \neq k} x_j x_k \right) = \lambda \left( \sum_{j \neq k} x_j x_k - \sum_{j=0}^{N-1} x_j + 1 \right).
\end{align}
In standard QUBO notation, off-diagonal interactions are summed over ordered pairs $j < k$, where $\sum_{j \neq k} x_j x_k = 2 \sum_{0 \leq j < k \leq N-1} x_j x_k$. Thus, the penalty contribution becomes:
\begin{align}
  P(\mathbf{x}) = \lambda \left( 2 \sum_{0 \leq j < k \leq N-1} x_j x_k - \sum_{j=0}^{N-1} x_j + 1 \right).
\end{align}

\textbf{Step 4 --- Combined Hamiltonian.} Adding the objective and penalty terms gives the full QUBO Hamiltonian:
\begin{align}
  H(\mathbf{x}) = \sum_{j=0}^{N-1} s_j x_j + \lambda \left( 2 \sum_{0 \leq j < k \leq N-1} x_j x_k - \sum_{j=0}^{N-1} x_j + 1 \right).
\end{align}

\textbf{Step 5 --- Q Matrix Entries and Offset.} Collecting coefficients for linear and quadratic terms yields the general closed-form expressions for the QUBO matrix $Q$ and constant offset $c$:
\begin{align}
  Q_{jj} &= s_j - \lambda, \quad \forall j \in \{0, \dots, N-1\}, \\
  Q_{jk} &= 2\lambda, \quad \forall 0 \leq j < k \leq N-1, \\
  c &= \lambda.
\end{align}

\paragraph{Penalty Weight Justification.}
The penalty weight $\lambda$ must be chosen such that any violation of the exactly-one constraint incurs an energy cost strictly greater than the maximum possible reduction in the objective function. Since the objective is a linear combination of scores $s_j$, the maximum gain from selecting an infeasible configuration (e.g., selecting zero or multiple candidates) is bounded by $\max_j |s_j|$. Therefore, setting $\lambda > \max_j |s_j|$ guarantees that infeasible states are energetically unfavorable. In practice, a sufficiently large constant such as $\lambda = 10000$ is employed to dominate typical hinge-loss magnitudes, ensuring constraint satisfaction without altering the relative ordering of feasible solutions.

\paragraph{Correctness Argument.}
Minimizing $H(\mathbf{x})$ over $\{0,1\}^N$ recovers the optimal candidate selection because feasible solutions satisfy $\sum_{j=0}^{N-1} x_j = 1$, rendering the penalty term identically zero. Consequently, the Hamiltonian reduces to the original objective function over the feasible subspace. For infeasible solutions, the penalty term evaluates to at least $\lambda$, which, by construction, exceeds any potential decrease in the objective term. Thus, the global minimum of $H(\mathbf{x})$ must reside in the feasible region and correspond to the candidate with the lowest regularized hinge-loss score.

\paragraph{Illustrative Example.}
Consider an instance with $N=3$ candidates and precomputed scores $s_0 = 12.5$, $s_1 = 8.3$, and $s_2 = 15.7$. Setting $\lambda = 10000$, the concrete objective is $12.5 x_0 + 8.3 x_1 + 15.7 x_2$. The concrete penalty terms expand to $10000 \left( 2x_0 x_1 + 2x_0 x_2 + 2x_1 x_2 - x_0 - x_1 - x_2 + 1 \right)$. The resulting Hamiltonian is:
\begin{align}
  H(\mathbf{x}) = 12.5 x_0 + 8.3 x_1 + 15.7 x_2 + 10000 \left( 2x_0 x_1 + 2x_0 x_2 + 2x_1 x_2 - x_0 - x_1 - x_2 + 1 \right).
\end{align}
The corresponding QUBO matrix entries and offset are:
\begin{align}
  Q_{00} &= -9987.5, & Q_{11} &= -9991.7, & Q_{22} &= -9984.3, \\
  Q_{01} &= 20000, & Q_{02} &= 20000, & Q_{12} &= 20000, \\
  c &= 10000.
\end{align}
Evaluating the three feasible configurations yields energies $H(1,0,0) = -9975.0$, $H(0,1,0) = -9983.2$, and $H(0,0,1) = -9968.6$. The global minimum occurs at $\mathbf{x}^* = (0, 1, 0)$ with energy $-9983.2$, correctly identifying candidate 1 as the optimal selection.

\end{document}